\section{Portability Across Metrics}
\label{sec:results-cross-metric}

\S\ref{sec:reversals} and \S\ref{sec:results-robustness} ask whether a
policy ranking under one fixed metric (ANWG) is portable across sources,
regions, and windows. A different question -- whether a ranking is
portable across \emph{metrics}, for the same source and region -- is
addressed here as an explicitly labeled post-campaign extension that
reuses the primary analysis's own statistical methods without modifying,
re-running, or re-deriving the underlying outcome matrix. Across all $3$
sources $\times$ $6$ regions $\times$ $55$
metric pairs among the study's tracked outcome metrics ($990$ conditions),
rank correlation between metrics is on average moderate (median Kendall's
$\tau_b = 0.419$, range $[-0.506, 1.000]$; $202/990$, $20\%$, of conditions
have \emph{negative} $\tau_b$), and the single top-ranked policy agrees
between the two metrics in only $316/990$ ($31.9\%$) of conditions
(disagreeing in the remaining $68.1\%$). At the level of individual
pairwise policy comparisons ($54{,}450$ total, one per (metric pair, condition, policy pair), generated from the 55 unique pairs among the 11 eligible tracked metrics: ANWG, completion fraction, weighted completion fraction, SLO-violation rate, weighted goodput, mean latency, p95 latency, request throughput, token throughput, mean TTFT, and p95 TTFT; p99 latency was ineligible and excluded), the same five-class taxonomy as \S\ref{sec:reversals} (ordinary, not
FDR-adjusted, bootstrap-CI classification; \S\ref{sec:uncertainty}) finds $3{,}244$ ($6.0\%$) show a statistically
supported practical disagreement in which
policy is preferred depending on which metric is used, versus $45{,}800$
($84.1\%$) in the same order, $5{,}205$ ($9.6\%$) a sign change too small
to be practically meaningful, and $201$ ($0.4\%$) an unsupported sign
change with a wide confidence interval; no aggregate comparison in this policy panel
was undefined or unestimable; individual policy/window observations can nevertheless be missing. Unlike \S\ref{sec:reversals}'s primary
reversal check, the separate BH-FDR correction removes some of these
cross-metric disagreements: of the $3{,}244$ raw-CI-supported records,
$3{,}207$ ($98.9\%$) also survive the within-family BH-FDR check, while
$37$ ($1.1\%$) do not, meaning the multiplicity-corrected count of
cross-metric disagreements is $3{,}207$ rather than $3{,}244$ -- a genuine,
disclosed sensitivity to the multiplicity correction, not an error in
either figure. Scheduler-ranking portability is
therefore not only source- and region-dependent (\S\ref{sec:results});
which metric is used to define ``better'' materially changes the ranking
in a non-trivial minority of conditions, and the median condition shows no
default expectation that the two metrics' top choices should agree.


\subsection{What a preferred-policy change means}
All 55 unordered pairs among 11 metrics are included at each of three
sources and six regions. Thus 990 denotes metric-pair conditions here;
it denotes policy-pair/source-pair/region records in the primary reversal
analysis. These are different denominators even though their sizes match.
The 54,450 records further expand each metric-pair condition into 55
policy-pair comparisons. They share windows, policies and metrics, so they
must not be interpreted as 54,450 independent replications.

Lower-is-better metrics are direction-normalized before comparison.
Kendall tau-b accommodates equal aggregate values; top-1 uses a stable
policy-ID tie break. Among the 1,980 metric sides of the 990 records,
1,010 have more than one policy tied at the maximum. The 31.9\% top-choice
agreement is exactly reproducible under the historical rule but should
not be described as disagreement between unique optimal policies. A
set-valued optimum-overlap analysis would be a separately specified
extension; it is not retroactively substituted for the reported statistic.

The median correlation describes this finite set of metric comparisons,
not a sampling distribution over all conceivable service objectives.
Correlated metrics, conditional observation sets and metric-specific
priority weighting can explain some disagreement by construction. The
analysis does not identify a uniquely correct deployment metric. It
requires a reader to state the intended objective before transporting a
ranking. Its post-campaign status limits confirmatory interpretation.
